\section{Implementation and Engineering}\label{sec:impl}

\subsection{Safety Footprint}
RedlineDB is written in safe Rust everywhere except where the C ABI
requires otherwise. The \texttt{redlinedb-ffi} crate contains
\texttt{unsafe} blocks because each entry point traffics in raw C
pointers and must dereference caller-supplied buffers; this is the
intended use of \texttt{unsafe} per the Rust safety
guidelines~\cite{rust_unsafe_guidelines,jung2021rustbelt}. Outside
\texttt{ffi}, the only \texttt{unsafe} is one well-audited
thread-local in \texttt{crates/kernel/src/engine/mod.rs} that arms
the per-thread fail-point context for the deterministic crash
harness; that block exists so a fail-point hit can panic the
specific worker thread without racing the rest of the
process. Everything else---kernel, SQL, public facade, bench
harness---compiles under \texttt{\#![forbid(unsafe\_code)]}-equivalent
review and depends on no \texttt{unsafe} idioms in user code paths.

\subsection{Lines of Code}
Table~\ref{tab:loc} reports the active source breakdown by crate. The
counts come from \texttt{find crates -name '*.rs' -not -path
'*/tests/*'} in the repository at \texttt{phase10-wave2-fused}. The
total active source is 47{,}565 lines (a +12{,}566 increase over
phase-9 driven primarily by the kernel-side integrity / JSON / vector /
HNSW / DiskANN modules and the SQL-side JSON1, vectorized executor,
collation, datetime, and REGEXP additions). For comparison, SQLite
ships an amalgamation of approximately 250{,}000 lines of
C~\cite{sqlite_footprint}.

\begin{table}[t]
\centering
\caption{Active Rust source LOC at \texttt{phase10-wave2-fused} vs SQLite reference~\cite{sqlite_footprint}}
\label{tab:loc}
\renewcommand{\arraystretch}{1.05}
\begin{tabular}{@{}lrr@{}}
\toprule
\textbf{Component} & \textbf{phase-9} & \textbf{phase-10} \\
\midrule
\texttt{redlinedb-kernel} (storage, WAL, MVCC, integrity, JSONB, vector) & 12{,}883 & 20{,}365 \\
\texttt{redlinedb-sql} (parser, planner, vectorized executor, JSON1) & 11{,}615 & 16{,}850 \\
\texttt{redlinedb} (public Rust facade) & 2{,}975 & 2{,}977 \\
\texttt{redlinedb-ffi} (C ABI shim) & 1{,}478 & 1{,}890 \\
\texttt{redlinedb-bench} (benchmark harness) & 5{,}144 & 5{,}483 \\
\midrule
\textbf{RedlineDB total active source} & \textbf{34{,}999} & \textbf{47{,}565} \\
\midrule
SQLite amalgamation (C) reference & \multicolumn{2}{r}{$\approx$250{,}000} \\
\bottomrule
\end{tabular}
\end{table}

\subsection{Failpoint Discipline}
The failpoint layer is a custom thin wrapper over the
\texttt{fail} 0.5 grammar~\cite{failrs} that retains
compile-neutrality when the feature is disabled. There are 16
failpoint sites threaded through WAL write, WAL fdatasync,
commit-publish, heap-mut, index-mut, catalog rename, and
checkpoint paths. The macro is shaped so that a no-op build
(release mode without the failpoint feature) compiles to no
instructions at the call site:

\begin{lstlisting}
let written = self.written_lsn;
let _ = written;
crate::fail_point!("wal::flush", |_| { Ok(written) });
self.active_file.sync_data()?;
\end{lstlisting}

A \texttt{failpoints::cfg} call validates the schedule grammar at
test setup, and a \texttt{cfg\_skip\_then\_panic} mode lets a test
skip $N$ hits and panic on the $(N+1)$-th to model
\texttt{kill\_after\_n\_hits} crash injection. The child-process
oracle records each acked commit to an fsync'd ack log; on restart,
recovery replays the WAL and the oracle confirms that every acked
commit is present and no ghost commits appear. Across the 24-case
failpoint matrix we observed 0 lost-acked-commits.

\subsection{Per-Transaction Index Undo}
The executor's per-tx undo log is a small enum:

\begin{lstlisting}
pub enum IndexUndoOp {
    Insert {
        index_id: IndexId,
        key: Vec<u8>,
        row: IndexRowRef,
    },
    DeleteMark {
        index_id: IndexId,
        key: Vec<u8>,
        row: IndexRowRef,
    },
}
\end{lstlisting}

On rollback (or on commit failure after kernel mutation) the executor
replays the inverse of each entry against the live index handle:
\texttt{Insert} is undone by \texttt{delete\_mark\_tx},
\texttt{DeleteMark} by \texttt{undelete\_mark\_tx}. Without this log,
a rolled-back \texttt{DELETE} would hide committed rows from indexed
reads, and a rolled-back \texttt{INSERT} could leak a stale
unique-conflict witness. The undo log lives in transaction-local
memory and is dropped on commit success; a long-running write
transaction therefore pays a per-mutation memory cost proportional
to its index-touch count, which we have not found to dominate any
realistic workload.

\subsection{Planner Gate}
The planner's index-access matcher refuses to advertise an index path
unless both the catalog entry and the engine handle agree the index
is consumable:

\begin{lstlisting}
for index in &table.indexes {
    if index.meta_page_id.is_none()
        || engine.index_handle(index.index_id).is_none() {
        continue;
    }
    // ... try to bind the leading key column ...
}
\end{lstlisting}

This is the contract that backs the \texttt{EXPLAIN} honesty story:
the executor and planner consume the same predicate for index
consumability, so a plan that says \texttt{IndexPointLookup} cannot
silently fall back to a table scan at runtime.

\subsection{Bench Harness}
The certification harness is a bin-packed parallel scheduler. Each
``cell'' (engine $\times$ workload $\times$ durability $\times$ thread
count $\times$ rep) is one child process; the scheduler packs cells
onto the host while reserving 4 cores for OS work and the reaper
thread, so a spike in another tenant cannot poison a measured run.
Every cell discards a real warmup run (the harness clears OS page
cache and walks a one-rep warmup; the warmup row is dropped from the
measured set). On Linux, an opt-in \texttt{--with-strace} flag
reroutes the child through \texttt{strace} so per-syscall counts can
be compared directly with the in-process pwrite/fdatasync counters.
Per-run telemetry includes throughput, p50/p95/p99/p999/max latency,
BUSY/LOCKED/timeout split error counts, RSS at peak,
fdatasync and pwrite counts, recursive on-disk \texttt{data\_bytes},
and an integrity checksum.

\subsection{Reproducibility}
Every certification run emits a \texttt{manifest.json} that pins the
git SHA, the Docker image digest, the host fingerprint (CPU model,
core count, kernel, file system), the SQLite PRAGMA snapshot used in
the benchmark, and per-output checksums. The repository tags each
certified run: \texttt{phase9-baseline}, \texttt{phase9-zlbap1-final},
\texttt{wave7-fused}, \texttt{phase9-xbabe1-certified}, and the
phase-10 markers \texttt{phase10-baseline},
\texttt{phase10-wave1-partial}, \texttt{phase10-wave2-fused},
\texttt{phase10-xbabe1-certified}, \texttt{phase10-fusion-green} are
the externally-visible markers in the public history. The PRAGMA
snapshot is recorded so that a reader who reproduces the run can
verify they are exercising SQLite under the same durability settings
we did (\texttt{journal\_mode=WAL}, \texttt{synchronous=FULL} for
strict cells, \texttt{synchronous=NORMAL} for relaxed cells).

\subsection{Phase 10: Long-Range Capabilities}
The phase-10 closure adds several capabilities the paper's first
release deferred to future work. Each landed under the same
parallel-agent + integrator-fusion discipline as the earlier
proof lanes, with one branch per lane and a single integrator running
the full proof matrix at each fusion point.

\textbf{Index MVCC.} Index entries now carry per-entry
\texttt{(create\_tx, delete\_tx)} tags rather than a boolean dead flag.
Visibility is checked via \texttt{point\_lookup\_visible} and
\texttt{range\_scan\_visible}, which take both a
\texttt{ConcurrentTxStatus} and a \texttt{Snapshot} so concurrent
readers and writers see consistent index contents without an
SQL-side undo log. The format version was bumped to v2 with a
v1$\rightarrow$v2 migration on \texttt{Engine::open}.

\textbf{Commit outcome.} \texttt{CommitOutcome::MaybeCommitted} flows
through the engine and the SQL executor so a post-fsync failure can
be surfaced to the caller as ``maybe committed,'' rather than as an
ordinary rollback that the application might retry against a state
that has already become durable.

\textbf{Heap/index integrity checker.}
\texttt{kernel/src/integrity/} provides a full equivalence check:
visible-row heap walk, full index tree dump, heap$\leftrightarrow$index
cross-check, page-checksum verification, and LSN monotonicity audit.
The result is a structured \texttt{IntegrityReport} per relation,
exposed via \texttt{PRAGMA redline\_index\_check} and
\texttt{PRAGMA redline\_full\_check}.

\textbf{Group-commit telemetry.} The WAL writer thread already
batched commits across waiters; phase-10 adds a 16-bucket batch-size
histogram (p50/p95/p99/max) on \texttt{WalSyncCounters}, an opt-in
per-core lane coordinator (default 1 lane), and a stub for a semantic
counter combiner gated behind \texttt{WalConfig::semantic\_combiner}.

\textbf{Vectorized executor and spillable sort.}
\texttt{sql/src/exec/vec/} adds selection vectors, a top-K min-heap
(emitted by the planner whenever \texttt{ORDER BY $\dots$ LIMIT $k$}
with $k\le 64$), hash aggregation with a spill lane, and an external
merge-sort that switches to disk runs when
\texttt{QueryMemoryBroker::work\_mem\_bytes} is exceeded. The new
bench workload \texttt{large-sort-spill} forces the spill path so
the harness validates the memory ceiling end to end.

\textbf{JSON.} The kernel ships a binary JSONB format
(\texttt{magic 0x96}, format-v1, type tags 0x00..0x08, LEB128
varints, zig-zag i64 encoding) plus a compiled JSON-path bytecode
with SIMD key compare. The SQL surface delivers SQLite JSON1
compatibility for \texttt{json}, \texttt{json\_array},
\texttt{json\_array\_length}, \texttt{json\_object},
\texttt{json\_extract}, \texttt{json\_set}, \texttt{json\_insert},
\texttt{json\_replace}, \texttt{json\_remove}, \texttt{json\_patch}
(RFC 7396), \texttt{json\_type}, \texttt{json\_valid},
\texttt{json\_quote}, \texttt{json\_minify}, the \texttt{->}
(JSON-text) and \texttt{->>} (SQL-value) operators, and SQLite's
shorthand path syntax (bare ident $\rightarrow$ \texttt{\$.ident},
bare integer $\rightarrow$ \texttt{\$[N]}). A 1000-iteration
round-trip fuzz validates the wire format against
\texttt{serde\_json::Value} ground truth.

\textbf{Vector search.} A native VECTOR column carries
dimension-checked \texttt{f32} arrays through the existing Blob
storage class (no new \texttt{SqlValue} variant). Distance kernels
expose L2, Cosine, and InnerProduct via runtime-dispatched
AVX2/NEON/scalar SIMD. The exact flat-scan operator computes top-$k$
in one pass; \texttt{<=>} (cosine) flows through the existing
binary-op path. An HNSW
index~\cite{malkov2018hnsw} (M=32, ef\_construction=200) achieves
recall@10 = 0.95 at ef\_search=64 on a 10k-vector 128-d Gaussian
synthetic dataset; a DiskANN-style Vamana
graph~\cite{subramanya2019diskann} (R=64, $\alpha$=1.2) achieves
recall@10 = 0.99 on a 10k-vector 32-d synthetic dataset, with the
4\,KiB sector layout designed in for future mmap-resident search.

\textbf{SQLite surface.} SAVEPOINT / RELEASE / ROLLBACK TO are
implemented via journal-and-replay on the SQL layer (no kernel-side
partial rollback hook is needed for the supported test matrix).
Multi-statement parsing returns the unconsumed remainder so the FFI
\texttt{sqlite3\_prepare\_v2} can populate \texttt{pzTail}; the
\texttt{sqlite3\_exec} path walks all statements and stops at the
first failing one with an \texttt{errmsg} owned by \texttt{CString}
and freed via \texttt{sqlite3\_free}. The conflict matrix
(\texttt{INSERT OR ABORT/FAIL/IGNORE/REPLACE/ROLLBACK}) routes through
a single \texttt{apply\_conflict\_resolution} helper across NOT NULL,
CHECK, UNIQUE, and PK constraints. Wrong-result fixes against
SQLite's documented behavior cover \texttt{SELECT ALL},
\texttt{NOT IN NULL}, NULL concatenation, divide / modulo by zero
returning NULL, scalar-function NULL propagation, CAST
truncation/prefix-parse, GLOB bracket / range / negation, and
ORDER BY honoring keys after GROUP BY / DISTINCT. New full-execution
features include REGEXP, the SQLite date/time function family
(\texttt{date}, \texttt{time}, \texttt{datetime}, \texttt{julianday},
\texttt{strftime}, \texttt{unixepoch} + modifiers), and the BINARY /
NOCASE / RTRIM collations. Foreign keys, ALTER TABLE DROP COLUMN,
partial / expression indexes, CTEs, CREATE VIEW, CREATE TRIGGER,
window functions, and generated columns parse cleanly and surface
structured ``not yet implemented'' errors at execute time so the
SQLite-compatibility surface is honest about what executes versus
what merely parses.

\textbf{Bench checksums.} Phase-9 used \texttt{MAX(k)} and
\texttt{COUNT(*)} as placeholder checksums, which an external audit
correctly flagged as not actually checksums. Phase-10 replaces them
with \texttt{DatasetChecksum}: \texttt{(row\_count, key\_xor,
payload\_hash)} folded across all rows, where \texttt{row\_hash} is
a deterministic XXH3-style fold of the row bytes and \texttt{key\_xor}
is order-independent. \texttt{DatasetChecksum::merge} is associative
so two re-runs of the same workload at the same seed produce
identical manifest checksums.

The phase-10 test matrix runs to 691 passing tests (vs.\ 241 at the
phase-9 wave-7 baseline; +450 phase-10 tests).
